\section{论文的主要内容和技术路线}

\subsection{主要研究内容}

本研究聚焦科学场景智能体工具调用的核心痛点，设置 4 项核心研究内容，具体如下：

\subsubsection{面向科学场景的工具抽象与统一接口规范设计}

本部分旨在搭建标准化科学工具封装体系，针对性解决当前科研领域工具接口不统一、跨工具集成难度大、调用与执行规范杂乱等问题，通过系统化架构设计与模块化落地，实现各类科研工具的规范化整合、通用化调用与安全化运行，为科研系统的工具联动、高效研发筑牢底层技术支撑，整套体系形成规范统一、模块封装、安全管控闭环衔接的完整架构。

为从根源统一工具调用标准，首先要设计通用科学工具抽象基类，明确规范工具名称、功能描述、参数格式、执行逻辑、结果返回形式及异常处理规则等核心要素，打破不同领域、不同类型科研工具的接口壁垒，让各类工具调用一致接口，减少接口差异带来的集成障碍，降低二次开发与跨工具适配成本\cite{schick2023toolformer,langchain2023}。

基于 Python 科学计算生态，聚焦科研场景封装核心工具模块，搭建可扩展的科学工具库，覆盖四大核心类别：数值计算类基于 NumPy、SciPy\cite{harris2020numpy,virtanen2020scipy}，封装线性代数运算、积分微分、函数拟合等基础计算工具；统计分析类基于 Pandas、SciPy\cite{harris2020numpy,mckinney2010pandas}，整合数据清洗、描述性统计、假设检验、回归分析等全流程统计工具；数据可视化类基于 Matplotlib，标准化封装折线图、散点图、热力图等科研常用绘图工具；四大模块既可独立使用，也能协同完成复杂科研任务。

搭建工具白名单机制与参数前置校验模块，拦截大模型生成的非法工具调用请求，校验输入参数合规性，过滤不符合规范的参数内容，从源头规避工具执行异常，保障系统运行安全性与稳定性。

\subsubsection{科学任务理解与拆解机制优化}

本部分旨在解决当前智能体处理专业科研任务时，普遍存在的任务拆解逻辑混乱、关键步骤遗漏、配套工具匹配不当、执行路径无法落地等痛点，全面提升智能体对科研场景的适配能力与任务处理精准度，为智能科研系统的高效调度执行筑牢核心支撑\cite{yao2023react,wang2023plansolve}。

搭建多轮迭代的层级化任务拆解框架，针对复杂科研任务实现"主任务 $\to$ 细分原子任务 $\to$ 可执行单步操作"的三级精细化拆解，确保每一个最小执行单元都能直接对接前文构建的标准化科学工具库，支持单个或组合工具调用，规避无对应工具支撑、无法落地执行的无效拆解步骤，强化任务拆解的可操作性\cite{yao2023react,shen2023hugginggpt}。

结合科研领域的专业特性深化 Prompt 工程优化\cite{wang2023plansolve}，融入多类型典型科研任务实操案例、细分领域基础科研规则与通用行业规范，进一步提升大模型对专业科研需求的语义理解与场景识别准确率，降低任务拆解错误率与工具误匹配率，使智能体能够稳定适配基础科研与专业细分领域的各类任务处理需求\cite{openai2023gpt4,wang2023plansolve}。

\subsubsection{工具调用优化与容错校验机制设计}

本部分聚焦科研场景下工具调用全流程优化，搭建容错与校验体系，核心目标是提升工具调用成功率、执行效率与结果可靠性，针对性解决当前科学工具调用容错性差、跨工具协同不畅、无效调用冗余、输出结果缺乏严谨性核验的突出问题\cite{schick2023toolformer,qin2023toollearning,shen2023hugginggpt}。

基于启发式规则构建动态工具选择与组合策略，基于任务拆解后的细化执行需求，综合考量工具功能适配度、执行效率与调用成本，匹配最优工具组合及标准化调用顺序，主动筛除重复冗余的调用步骤，压缩无效执行环节，减少资源损耗，提升整体任务执行效率\cite{yao2023react,qin2023toollearning}。

搭建跨工具统一数据流管理模块，制定中间数据格式规范，打通不同工具间的数据传输壁垒，实现前序工具输出结果与后序工具输入参数的自动格式转换、无缝适配，解决多工具协同执行中的数据不兼容、格式不匹配等问题，保障链式工具调用的连贯性与流畅性\cite{shen2023hugginggpt,langchain2023}。

搭建覆盖全调用流程的容错处理与结果双重校验机制，一方面构建自动化异常处理闭环，针对调用失败、参数报错、执行中断等问题，完成错误原因自动溯源、参数智能修正、按需重试执行，避免单次异常导致整体任务中断\cite{schick2023toolformer,qin2023toollearning}；另一方面嵌入专业科研校验规则，封装统计显著性检验、数据合理性筛查、结果逻辑验证等专项工具，对输出结果开展自动化复核，严格把控结果质量，确保所有输出内容完全符合科研工作的严谨性与规范性要求\cite{jumper2021alphafold,butler2018ml4mat}。

\subsubsection{系统原型实现与场景验证}

本部分是研究的落地验证环节，旨在验证系统的可用性与优化效果。

完成各模块的集成开发，构建完整的系统原型，实现用户科研需求输入 $\to$ 任务自动拆解 $\to$ 工具调用执行 $\to$ 结果校验输出的全流程自动化运行，并完成系统的全流程联调与 bug 修复。

构建科学任务测试数据集，涵盖 20 个不同类型、不同难度的典型科研任务，设计对比实验，以任务成功率、工具调用准确率、平均执行步数为核心指标，对比本研究优化后的系统与原生 GPT-4 直接调用\cite{openai2023gpt4}、通用 LangChain Agent\cite{langchain2023} 的性能差异，验证优化方案的有效性。

选取真实科研场景，包括实验数据统计与可视化、数值计算与拟合分析等，完成系统的场景化验证，测试系统在真实科研场景中的可用性与实用性。

\subsection{技术路线}

本研究基于 Python 技术栈与主流开源框架完成系统开发。整体技术路线遵循循序渐进、闭环迭代的原则，划分为 6 大核心实施环节，具体实施流程如下。

\subsubsection{前期筹备与文献调研梳理}

本环节完成领域文献系统性调研与技术选型落地。

梳理大语言模型工具调用、智能体框架构建及 AI for Science 领域国内外核心文献、顶会顶刊研究成果与优质开源项目\cite{brown2020gpt3,yao2023react,schick2023toolformer,jumper2021alphafold}，总结现有研究的优势亮点与现存短板，完成开题报告的撰写、修订与审核。完成全套技术选型，初步确定适配科研场景的开发环境与技术栈体系：底层大模型基于通义千问 API 实现核心推理能力，备选 OpenAI GPT API\cite{openai2023gpt4}；科学计算模块基于 NumPy、SciPy、Pandas、Matplotlib 等模块\cite{harris2020numpy,virtanen2020scipy,mckinney2010pandas}。

\subsubsection{标准化科学工具库封装实现}

本环节完成核心科学工具库的落地开发与测试。

完成数值计算、统计分析、数据可视化三大类高频科研工具的封装，并分别针对每一项独立工具编写各自对应的单元测试代码，验证工具功能的正确性与稳定性。

完成工具白名单管控模块与参数前置校验模块开发，实现非法调用与异常参数的前置拦截\cite{qin2023toollearning}，完成工具库整体联调测试，确保整套工具库满足后续系统调用的稳定性需求。

\subsubsection{科学任务拆解机制优化开发}

本环节开展科研任务拆解模块的优化与开发工作。

设计适配科学任务的思维链 Prompt 模板与层级化任务拆解模板\cite{wang2023plansolve}，搭建标准化任务拆解流程，完成模块核心代码开发。构建专项测试任务集，对 Prompt 模板进行多轮调试与迭代优化，重点提升智能体对科研任务的理解准确率、拆解逻辑严谨性，减少无效拆解与逻辑漏洞，保障任务拆解结果可落地、可执行\cite{openai2023gpt4,yao2023react}。

\subsubsection{工具调用与容错校验模块研发}

本环节完成动态工具选择与组合策略、跨工具数据流管理模块、自动化异常重试机制、科研结果合规性校验模块的开发工作。

针对每个独立模块开展专项功能测试与异常场景测试，验证动态工具匹配效率、跨工具数据适配能力、异常处理响应速度与结果校验精准度\cite{schick2023toolformer,qin2023toollearning,shen2023hugginggpt}，确保各模块功能达标、运行稳定，为系统全流程流畅执行提供保障。

\subsubsection{系统集成与实验验证分析}

本环节整合集成前期开发的所有独立模块，搭建完整的系统原型，开展全流程联调测试，排查并修复运行过程中出现的漏洞与 bug，优化系统运行流畅度与交互可用性\cite{langchain2023,llamaindex2023}。

构建涵盖多类科研场景的标准化测试数据集，设计对照实验与性能验证实验，采集并整理实验数据，量化分析系统优化效果。选取真实科研场景开展落地验证\cite{jumper2021alphafold,butler2018ml4mat}，根据测试反馈进一步优化系统适配性与实用性，确保系统满足实际科研工作需求。

\subsubsection{论文撰写与最终成果定稿}

完成毕业论文（设计）全文撰写，涵盖研究背景与意义、国内外相关工作综述、系统整体设计、核心模块实现、实验设计与结果分析、研究结论与展望等核心章节。结合导师指导意见，对论文内容进行多轮修订完善，同步补充优化系统源码、实验数据、测试报告等配套毕设材料，最终完成论文定稿与全部毕业设计材料的整理归档。

\subsection{可行性分析}

\subsubsection{技术可行性}

本研究所有核心功能基于成熟的开源框架与技术栈实现。底层大模型采用成熟商用 API，具备稳定的函数调用能力\cite{openai2023gpt4}；科学工具基于 Python 科学计算开源库封装\cite{harris2020numpy,virtanen2020scipy,mckinney2010pandas}，这些库经过长期的工业界与学术界验证，功能稳定、文档齐全。整体技术方案成熟可靠，技术可行性高。

\subsubsection{资源可行性}

本研究底层大模型采用商用 API，通过按量付费的方式使用；相关的文献资料、开发教程、开源案例等可通过知网、IEEE Xplore、GitHub、官方文档等公开渠道获取\cite{langchain2023,llamaindex2023}。资源可行性高。

\subsubsection{工作量与实现难度可行性}

本研究核心开发工作包括工具封装、Prompt 工程、规则模块开发、系统集成等，预计可在 12 周内完成核心系统的开发与测试；实验环节需构建测试数据集、完成对比实验，无需大量的算力消耗。同时，本人具备较为丰富的 Python 开发、科学计算库使用、大模型 API 调用等相关经验，能够胜任系统开发工作。工作量与实现难度符合要求。
